\SourceLeaf{35}
\fancyhead[L]{\small General principles of the theory of the shape of fluids}
\renewcommand{\Article}[1]{\par\medskip\Needspace{4\baselineskip}\pdfbookmark[1]{Art. #1}{fluidorum-art-#1}\centerline{\textbf{#1.}}\nobreak\smallskip}
\pdfbookmark[0]{General Principles of the Theory of the Shape of Fluids in Equilibrium}{fluidorum-body}
\par\medskip
\centerline{GENERAL PRINCIPLES}
\centerline{OF THE THEORY OF THE SHAPE OF FLUIDS}
\centerline{IN EQUILIBRIUM.}
\medskip
\expandafter\def\csname HistoricalNoteText4\endcsname{It is established that the greatest attraction a given homogeneous mass can exert on a given point under that law is to the attraction the same mass, formed into a sphere, exerts on a point on its surface as $3$ is to $\BareRoot[3]{25}$. The latter attraction, in turn, is easily compared with gravity.}
\expandafter\def\csname HistoricalNoteText5\endcsname{This is how most of the objections raised in the Pavia journal (Giornale di fisica etc. T. 9) should be judged. The distinguished \textsc{Petit} answered them aptly in Annales de chimie et de physique T. 4.}
The perceptive \textsc{Clairaut} was the first to enumerate with insight and precision the forces governing the rise or depression of fluids in capillary tubes. But he left the law of the forces entirely untouched, so his enumeration could yield nothing toward a mathematical explanation of the phenomena. Ordinary attraction, inversely proportional to the square of the distance, explains all celestial motions with such success, yet it is of no use in explaining capillary phenomena or those of adhesion and cohesion. A correctly carried out calculation readily shows that, under this law, the attraction of any body with which experiments can be performed, that is, any body whose bulk can be treated as nothing in comparison with the entire Earth, vanishes in comparison with gravity at a point wherever it lies, even if it is in contact with the body\HistoricalNote{4}. The proper conclusion is that at the smallest distances this law of attraction no longer agrees with nature and requires some modification. Equivalently, the particles of bodies exert another force alongside that attraction, noticeable only at the smallest distances. All the phenomena point to this second part of the attractive force (\textit{mole\SourceLeaf{36}cular attraction}) being imperceptible even at the smallest distances we can measure, while at imperceptibly small distances it can far exceed the first part, which is inversely proportional to the square of the distance.

The distinguished \textsc{Laplace} proceeded from this single assumption about the nature of molecular forces, leaving the law of their decrease with increasing distance otherwise entirely undetermined. He was the first to calculate their effect on the shape of fluid surfaces with precision. Having established a general equation for the equilibrium shape, he tried to use it to explain both capillary phenomena in the strict sense and many related phenomena. The remarkable agreement with careful experiments has confirmed these investigations throughout. They rank among the finest advances in natural philosophy that we owe to that great geometer. The objections some authors have raised against them are, for the most part, of little or no consequence\HistoricalNote{5}.

There are, to be sure, points in the calculations of the distinguished \textsc{Laplace} that do not fully accord with rigorous reasoning. In the earlier memoir, \textit{théorie de l’action capillaire}, $\varphi f$ denotes the strength of attraction at distance $f$, and the integral $\int\varphi f.\mathrm{d}f$ taken from $f=x$ to $f=\infty$ is set $=\Pi x$. Next the integral $\int\Pi f.f\mathrm{d}f$ taken from $f=x$ to $f=\infty$ is set $=\Psi x$. Finally the values of the integrals $2\pi\int\Psi f.\mathrm{d}f$, $2\pi\int\Psi f.f\mathrm{d}f$ taken from $f=0$ to $f=\infty$ are set respectively $=K$ and $=H$, with $\pi$ denoting half the circumference of a circle of radius $=1$. The nature of the function $\varphi f$ is left entirely unspecified, provided it is imperceptible for all perceptible values of $f$. From this assumption \textit{alone}, however, it would by no means follow that $\Pi f,\Psi f$ must also be imperceptible for perceptible values of $f$. Nor is there any more reason to infer, as the memoir states, that the integrals $2\pi\int\Psi f.\mathrm{d}f$, $2\pi\int\Psi f.f\mathrm{d}f$ taken from $f=0$ to a perceptible \textit{finite} value of $f$ differ imperceptibly from $K,H$. One could imagine infinitely many forms of the function $\varphi f$ satisfying the fundamental assumption for which these conclusions would be wrong. Indeed, if $\varphi f$ is taken to express the total attraction, it will also contain a part of the form $\frac{a}{ff}$, on which ordinary attraction depends. Even if this term is to be regarded as imperceptible when the di\SourceLeaf{37}mensions of attracting bodies such as those used in experiments are imperceptible in comparison with the entire Earth, the second integration, if extended to infinity, would already introduce an infinite term into the function $\Psi f$.

If these and similar points appear to suggest some slight carelessness, it certainly concerns the manner of exposition more than the substance. The second memoir, \textit{Supplément à la théorie de l’action capillaire}, makes it clear that the distinguished \textsc{Laplace} tacitly understood $\varphi f$ to mean the part added to ordinary attraction, rather than the total attraction. It is readily apparent that ordinary attraction can make no perceptible difference in our experiments. He even indicates that he thinks of the function $\varphi f$ as like the exponential $e^{-if}$, with $i$ denoting a very large quantity, or rather with $\frac{1}{i}$ denoting a very small length. There is no need, however, to restrict the generality so much. Anyone who looks to the substance more than the words can readily see that it suffices to extend these integrations only to an arbitrary perceptible distance instead of to infinity, or, if one prefers, to a finite distance greater than the dimensions encountered in experiments.

The theory has another and much more serious defect, one that, so far as we know, its critics did not even notice. It has two parts. One establishes a general equation between the partial differentials of the coordinates for the free surface of the fluid. This equation depends on the molecular force of attraction that the fluid particles exert on one another, and this part of the theory is so fully worked out that nothing essential remains to be desired. But an equation between partial differentials of this kind, whose integration, if analysis could perform it, would introduce arbitrary functions, is insufficient to determine the shape of the surface \textit{completely}. That requires a \textit{new} condition determining the nature of the shape at the boundaries and such a condition is supplied by the second part of the theory. It requires that the angle between the plane tangent to the free surface of the fluid where it meets the vessel, or more precisely at the limit of the perceptible attractive force of the vessel wall, and the plane tangent to the vessel wall at that same place be \textit{constant}. This angle is determined by the relation between the strengths of the molecular forces of the vessel and the fluid, provided that the shape of the vessel is continuous where it meets the free surface of the fluid. Yet the distinguished \textsc{Laplace} did not even undertake to demonstrate this central proposition of the entire theory by calculation. The relevant considerations pre\SourceLeaf{38}sented on p. 5 of the earlier memoir amount only to a vague argument and already assume what had to be demonstrated. The calculations undertaken on p. 44 and following yield no result. The second memoir does treat the rise of fluid in capillary tubes by another method, whose outcome, when compared with that of the earlier method, supplies a formula for that angle between the tangent planes. The formula is indeed correct. But it must be noted that the constancy of the angle is already \textit{assumed} here, strictly speaking. Moreover, the method, itself rather unsatisfactory, is restricted to the very special case of a prismatic vessel with vertical walls. In view of this, it must be admitted that the theory proposed by the distinguished \textsc{Laplace} still falls short and remains incomplete in essential respects.

We will therefore take up again from the beginning the theory of the equilibrium shape of fluids under the action of gravity and of the molecular forces of both the fluid itself and the vessel. We will follow an entirely different method, drawn from the first principles of dynamics, and adopt the utmost generality from the outset. This investigation will lead to a striking new theorem that brings the complete theory together in a single, exceedingly simple formula. Both parts of the distinguished \textsc{Laplace}’s theory will follow naturally from this theorem.
\par\medskip\centerline{\rule{.3\textwidth}{.3pt}}

